\subsection{正则局部环的定义}

令 A 为诺特局部环。设 \((x_{i})_{i\in I}\) 是 \(m_{A}\) 中元素的一个族。根据 II, § 3, n° 2 的命题 4 的推论 2，假设族 \((x_{i})_{i\in I}\) 生成 A 的理想 \(m_{A}\)，或假设 \(x_{i}\) 模 \(m_{A}^{2}\) 的类生成 \(\kappa_{A}\) -向量空间 \(m_{A}/m_{A}^{2}\)，两者是等价的；若如此，由 § 3, n° 2 的推论，有 dim(A) ≤ Card(I)。因而有不等式

\[
\dim (A) \leqslant [ m _ {A} / m _ {A} ^ {2}: \kappa_ {A} ] \leqslant \operatorname{Card} (I)\tag{1}
\]

对于每个族 \((x_{i})_{i\in I}\)，若其生成 A 的理想 \(m_{A}\)，则都有此不等式。

定义 1. --- 如果有 \(\dim(A) = [m_A/m_A^2:\kappa_A]\)，则称诺特局部环 A 为正则的。A 的坐标系是指 \(m_A\) 中的每个元素族，其模 \(m_A^2\) 的类构成 \(\kappa_A\)-向量空间 \(m_A/m_A^2\) 的一个基。

因此，正则诺特局部环 A 的一个坐标系是有限族 \((x_{i})_{i\in I}\)，它生成 A 的理想 \(m_{A}\)，并满足 \(\text{Card}(I) = \dim(A)\)。反之，若诺特局部环 A 的理想 \(m_{A}\) 由 d 个元素生成，且 \(d \leqslant \dim(A)\)，则环 A 是正则的。

例子。--- 1) 维数为 0（分别为 1）的正则诺特局部环是域（分别为离散赋值环）（VI, § 3, n° 6, 命题 9 以及下文定理 1 的推论 1）。设 A 是离散赋值环；则 m \(_{A}\) 中的元素 t 是一致化元，当且仅当 \{t\} 是 A 的一个坐标系。

\begin{enumerate}
\def\labelenumi{\arabic{enumi})}
\setcounter{enumi}{1}
\tightlist
\item
  令 \(k\) 为域，令 \(n\) 为正整数。形式幂级数环 \(k[[X_1, ..., X_n]]\) 是维数为 \(n\) 的正则诺特局部环（§ 3, n° 4, 命题 8 的推论 3）。令 \(F_1, ..., F_n\) 为 \(k[[X_1, ..., X_n]]\) 中无常数项的形式幂级数；序列 \((F_1, ..., F_n)\) 成为 \(k[[X_1, ..., X_n]]\) 的一个坐标系的充要条件是矩阵 \(\left(\frac{\partial F_i}{\partial X_j}(0, ..., 0)\right)\) 可逆。
\end{enumerate}

更一般地，设 A 是维数为 \(r\) 的正则诺特局部环；则 \(A[[X_1, ..., X_n]]\) 是维数为 \(r + n\) 的正则诺特局部环。若 \((a_1, ..., a_r)\) 是 A 的一个坐标系，则 \((a_1, ..., a_r, X_1, ..., X_n)\) 是 \(A[[X_1, ..., X_n]]\) 的一个坐标系。

\begin{enumerate}
\def\labelenumi{\arabic{enumi})}
\setcounter{enumi}{2}
\tightlist
\item
  令 A 是维数为 \(r\) 的完备正则诺特局部环。受限形式幂级数环 A \(\{X_1, ..., X_n\}\) 是维数为 \(r + n\) 的正则诺特局部环（§ 3, n° 4, 注 1）。若 \((a_1, ..., a_r)\) 是 A 的一个坐标系，则 \((a_1, ..., a_r, X_1, ..., X_n)\) 是 A \(\{X_1, ..., X_n\}\) 的一个坐标系。
\end{enumerate}

* 4) 令 \(k\) 为完备非离散赋值域。\(n\) 元形式幂级数在 \(k^n\) 中 0 的某邻域内收敛时所成的环，是维数为 \(n\) 的正则诺特局部环 ( \(\S\) 3, n° 4, 注 2). *

\begin{enumerate}
\def\labelenumi{\arabic{enumi})}
\setcounter{enumi}{4}
\tightlist
\item
  令 \(k\) 为域，A 为有限型 \(k\)-代数且为整环，\(\mathfrak{m}\) 为 A 的极大理想。诺特局部环 \(\mathrm{A}_{\mathfrak{m}}\) 是正则的，当且仅当有 \(\dim(\mathrm{A}) = [\mathfrak{m}/\mathfrak{m}^{2}:\mathrm{A}/\mathfrak{m}]\)：事实上，有 \(\dim(\mathrm{A}_{\mathfrak{m}}) = \dim(\mathrm{A})\) ( \(\S 2\) , \(no. 4\) , 定理 3 的推论 2)，并且关于域 A/m 的向量空间 \(\mathfrak{m}/\mathfrak{m}^{2}\) 与 \(\mathrm{mA}_{\mathfrak{m}}/(\mathrm{mA}_{\mathfrak{m}})^{2}\) 同构（II, \(\S 3\) , \(no. 3\) , 命题 9）。特别地，若 \(k\) 代数闭，所述条件等价于 \(\dim(\mathrm{A}) = [\mathfrak{m}/\mathfrak{m}^{2}:k]\)（V, \(\S 3\) , \(no. 3\) , 命题 1）。
\end{enumerate}

* 6) 设 X 是完美域 \(k\) 上的代数簇。X 在点 \(x\) 处非奇异，当且仅当 X 在 \(x\) 处的局部环是正则的。*

* 7) 设 A 是正则诺特局部环。稍后将看到，对于 A 的每个素理想 p，诺特局部环 \(A_{p}\) 都是正则的。 *

命题 1.--- 设 A、B 是诺特局部环，\(\rho:A\to B\) 是一个局部同态，并使 B 成为平坦 A-模。设有 \(m_{B} = B \cdot \rho(m_{A})\)。则有 \(\dim(A) = \dim(B)\)，且 B 正则当且仅当 A 正则。

第一条断言由 § 3, no. 4 的命题 7 的推论 1 得出。由于 B 在 A 上平坦，可将 \(\mathfrak{m}_{\mathrm{B}}^{k} = \mathbf{B}.\rho (\mathfrak{m}_{\mathrm{A}}^{k})\) 认作 \(\mathbf{B}\otimes_{\mathbf{A}}\mathfrak{m}_{\mathbf{A}}^{k}\)，对于每个满足 \(k\geqslant 0\) 的整数，进而将 \(\mathfrak{m}_{\mathrm{B}} / \mathfrak{m}_{\mathrm{B}}^{2}\) 认作 \(\mathbf{B}\otimes_{\mathbf{A}}(\mathfrak{m}_{\mathrm{A}} / \mathfrak{m}_{\mathrm{A}}^{2})\)，或也可认作 \(\kappa_{\mathrm{B}}\otimes_{\kappa_{\mathrm{A}}}(\mathfrak{m}_{\mathrm{A}} / \mathfrak{m}_{\mathrm{A}}^{2})\)。因而有

\[
\left[ \mathfrak {m} _ {\mathrm{B}} / \mathfrak {m} _ {\mathrm{B}} ^ {2}: \kappa_ {\mathrm{B}} \right] = \left[ \mathfrak {m} _ {\mathrm{A}} / \mathfrak {m} _ {\mathrm{A}} ^ {2}: \kappa_ {\mathrm{A}} \right],\tag{2}
\]

命题立即由此得到。

推论。--- 诺特局部环 A 正则，当且仅当其完备化 \(\hat{A}\) 正则。

事实上，\(\hat{\mathbf{A}}\) 在 \(\mathbf{A}\) 上平坦，且有 \(\mathfrak{m}_{\hat{\mathbf{A}}} = \hat{\mathbf{A}}. \mathfrak{m}_{\mathbf{A}}\)（III, § 3, no. 4, 定理 3 以及 § 2, no. 12, 命题 16 的推论 2）。

